% Keep this static CV consistent with the web CV and content collections.
% After editing either version, run the content tests and inspect the generated PDF.
\documentclass[11pt]{article}
\usepackage{cmap}
\usepackage[T1]{fontenc}
\usepackage[utf8]{inputenc}
\usepackage{charter}
\usepackage{microtype}
\DisableLigatures{encoding=*, family=*}
\usepackage[margin=0.75in]{geometry}
\geometry{a4paper}
\usepackage[hidelinks,
  pdftitle={Saheed Faremi - Curriculum Vitae},
  pdfauthor={Saheed Faremi},
  pdfsubject={Research and engineering curriculum vitae},
  pdfkeywords={EEG microstates, generative models, machine learning, data science}
]{hyperref}

\pagestyle{empty}
\setlength{\parindent}{0pt}
\setlength{\parskip}{0pt}
\linespread{0.9}
\hyphenpenalty=10000
\exhyphenpenalty=10000
\emergencystretch=1.5em

% Standard LaTeX list lengths, adjusted without a list-formatting package.
\setlength{\leftmargini}{1.35em}
\setlength{\leftmarginii}{1.1em}
\setlength{\labelwidth}{0.8em}
\setlength{\labelsep}{0.45em}
\setlength{\itemindent}{0pt}
\setlength{\listparindent}{0pt}
\setlength{\topsep}{1pt}
\setlength{\partopsep}{0pt}
\setlength{\parsep}{0pt}
\setlength{\itemsep}{1pt}

\newcommand{\cvsection}[1]{%
  \par\vspace{5pt}%
  {\fontsize{11.5pt}{13pt}\selectfont\bfseries\MakeUppercase{#1}\par}%
  \nobreak\vspace{2pt}\hrule height 0.4pt\nobreak\vspace{2pt}%
}
\newcommand{\cventry}[2]{\noindent\textbf{#1}\hfill #2\par}
\newcommand{\skill}[2]{\noindent\textbf{#1:} #2\par}

\begin{document}

\begin{center}
  {\fontsize{20pt}{23pt}\selectfont\textbf{Saheed Faremi}\par}
  \vspace{2pt}
  PhD researcher, University College Cork · Data scientist\par
  \vspace{2pt}
  \href{mailto:saheedfaremi@gmail.com}{saheedfaremi@gmail.com} · Available to work anywhere in Ireland\par
  \href{https://github.com/fayisode}{GitHub} ·
  \href{https://www.linkedin.com/in/faremi-saheed}{LinkedIn} ·
  \href{https://scholar.google.com/citations?user=tw1P0nsAAAAJ}{Google Scholar} ·
  \href{https://orcid.org/0009-0007-8428-7019}{ORCID}\par
\end{center}
\vspace{2pt}
PhD researcher and data scientist with five years of experience building production software. I study what brief brain states represent and how they change over time. In industry, I apply machine learning to healthcare, finance and compliance.\par
{\small Last updated November 2026\par}

\cvsection{Research areas}
\skill{Brain and behaviour}{EEG microstates and brain-signal analysis; neuroscientific knowledge discovery; neuroimaging; nonlinear time-series analysis}
\vspace{1.5pt}
\skill{Explainable and reproducible AI}{Explainable Artificial Intelligence; mechanistic interpretability; reproducible machine learning}
\vspace{1.5pt}
\skill{Applications}{Deep learning for health and behaviour; production AI and data systems}

\cvsection{Technical skills}
\skill{Brain-signal analysis}{PyTorch; scikit-learn; MNE}
\vspace{1.5pt}
\skill{Applied AI and agents}{Azure OpenAI; AI agents and tool use; prompt engineering; retrieval-augmented generation (RAG); document question answering; model fine-tuning}
\vspace{1.5pt}
\skill{Software \& cloud engineering}{Python; Go; TypeScript; SQL; FastAPI; Azure; CI/CD; SLURM; multi-GPU scheduling; memory-aware job allocation; automated testing}
\vspace{1.5pt}
\skill{Research evaluation}{Testing across participants; statistical significance tests; correction for multiple comparisons; time-series and duration analysis}

\cvsection{Education}
\cventry{PhD, University College Cork}{In progress}
Neural engineering and machine learning. Supervisor: Luca Longo.\par
\vspace{2pt}
\cventry{MSc Computer Science, Technological University Dublin}{Completed}
Machine learning, deep learning, data mining, visualisation and statistics. 4.0/4.0 GPA (80/100).\par
\vspace{2pt}
\cventry{BSc Information Technology, University of Eswatini}{Completed}
Machine learning, deep learning, data mining, visualisation and statistics. Second Class Honours, Upper Division (2.1).\par

\enlargethispage{2\baselineskip}
\cvsection{Experience}
\cventry{PhD researcher}{In progress}
\textbf{Artificial Intelligence and Cognitive Load Research Lab (AICL), University College Cork}\par
\begin{itemize}
  \item Compared image-based and electrode-based representations of brain activity on 203 recordings. The electrode-based approach explained 61.8\% of the variation, versus 54.2\% for the image-based approach, which grouped similar activity better.
  \item Analysed brain-activity recordings from 203 people to understand how activity changes from one state to another. Paired statistical tests showed that accounting for state duration represented these changes more accurately than the current state alone, improving the result from 55.7\% to 63.7\%.
  \item Showed that analysis settings can change average reported brain-state durations from 402 to 24 milliseconds. Using the same processing method for both approaches changed the image-based result from 143.2 to 71.0 milliseconds, close to the electrode result of 70.2 milliseconds.
  \item Evaluated 4,832 model configurations on SLURM-managed GPU and IBM Power9 systems, scoring 110 measures per model and testing consistency across participants.
\end{itemize}
\vspace{2pt}
\cventry{Data scientist to Founding engineer}{May 2021 to present}
\textbf{Etihuku (Curnance)}\par
\begin{itemize}
  \item Led a four-person team to develop a compliance dashboard for tracking legal obligations, document submissions, compliance scores and risk levels, reducing manual handling by 60\%.
  \item Built a facial-recognition system for Know Your Customer (KYC) checks, using identity documents and liveness verification to flag potentially fraudulent user profiles for review. Developed automated checks to identify suspicious financial transactions, improving fraud detection by up to 35\%.
  \item Built an LLM proposal generator from scratch for business bids, using RAG to turn past proposals, contracts and client requirements into seven-section drafts and cut manual preparation time by 80\%.
\end{itemize}

\newpage
{\small Saheed Faremi · Curriculum Vitae · Page 2 of 2\par}
\vspace{3pt}\hrule height 0.4pt

\cvsection{Publications}
\noindent\textbf{Peer-reviewed}\par
\vspace{2pt}
\begin{enumerate}
  \setlength{\itemsep}{2pt}
  \item \href{https://aics2026.com/}{Temporal Dynamics and Latent-Space Interpretability of Generative EEG Microstate Segmentation}.\newline
  \textbf{Saheed Faremi}, Andrea Visentin, Luca Longo.\newline
  \emph{AICS 2026, University College Dublin (accepted; proceedings to appear)}, 2026.

  \item \href{https://link.springer.com/article/10.1186/s40708-026-00327-9}{Integrating Convolutional Variational Autoencoders and the Gaussian Mixture Model for efficient manifold learning and clustering of spatially preserved EEG topographic maps}.\newline
  \textbf{Saheed Faremi}, Luca Longo.\newline
  \emph{Brain Informatics}, 2026. DOI: \href{https://doi.org/10.1186/s40708-026-00327-9}{10.1186/s40708-026-00327-9}.

  \item \href{https://doi.org/10.1109/ICTAS59620.2024.10507142}{Machine Learning Models for Identifying Factors Influencing and Predicting Malaria Among Children Under Five Years in Nigeria}.\newline
  \textbf{Akinpelumi Saheed Faremi}, Boluwaji Akinnuwesi, Elliot Mbunge, Petros M. Mashwama, Stephen Fashoto, Polite Zenzo Ncube, John Batani, Shamsudeen Ademola Sanni, Yinusa A. Faremi, Andile Metfula.\newline
  \emph{IEEE ICTAS 2024}, 2024. DOI: \href{https://doi.org/10.1109/ICTAS59620.2024.10507142}{10.1109/ICTAS59620.2024.10507142}.
\end{enumerate}
\vspace{5pt}
\noindent\textbf{Preprints, registered protocols \& doctoral work}\par
\vspace{2pt}
\begin{enumerate}
  \setcounter{enumi}{3}
  \setlength{\itemsep}{2pt}
  \item \href{https://arxiv.org/abs/2605.10947}{Interpretable EEG Microstate Discovery via Variational Deep Embedding: A Systematic Architecture Search with Multi-Quadrant Evaluation}.\newline
  \textbf{Saheed Faremi}, Andrea Visentin, Luca Longo.\newline
  \emph{XAI 2026, Late-breaking work and Doctoral Consortium track, Fortaleza, Brazil (arXiv preprint)}, 2026. DOI: \href{https://doi.org/10.48550/arXiv.2605.10947}{10.48550/arXiv.2605.10947}.

  \item \href{https://osf.io/v2w3z}{Autoencoder-Based Models for Scalp EEG: A Systematic Review of Architectures, Applications, Latent Representations, Interpretability, and Validation}.\newline
  Morteza Akbari, \textbf{Saheed Akinpelumi Faremi}, Luca Longo.\newline
  \emph{OSF Registries (registered systematic-review protocol)}, 2026. DOI: \href{https://doi.org/10.17605/osf.io/v2w3z}{10.17605/osf.io/v2w3z}.

  \item Explainable Disentangled Representation Learning of Recurring Brain Activation Patterns via Variational Autoencoders.\newline
  \textbf{Saheed Faremi}.\newline
  \emph{XAI World Conference 2025, Doctoral Proposals track}, 2025.
\end{enumerate}
\vspace{5pt}
\noindent\textbf{Under review and manuscripts}\par
\vspace{2pt}
\begin{enumerate}
  \setcounter{enumi}{6}
  \setlength{\itemsep}{2pt}
  \item Convolutional Variational Deep Embedding Produces Canonical EEG Microstate Maps and Richer Transition Dynamics than Modified K-Means on Resting EEG.\newline
  \textbf{Saheed Faremi}, Andrea Visentin, Luca Longo.\newline
  \emph{Cognitive Computation (manuscript in preparation)}, 2026.

  \item Spatial Encoding Geometry Determines the Maps and Transitions Recovered by Deep EEG Microstate Models.\newline
  \textbf{Saheed Faremi et al.}\newline
  \emph{ICLR 2027 (under review; OpenReview \#50003)}, 2026.

  \item Prediction of Early Detection of Autism Spectrum Disorder in Toddlers Using Explainable Artificial Intelligence.\newline
  \textbf{Stephen Fashoto et al.}\newline
  \emph{Under review}, 2026.
\end{enumerate}

\cvsection{Talks}
\begin{itemize}
  \setlength{\itemsep}{3pt}
  \item \href{https://aics2026.com/}{Temporal Dynamics and Latent-Space Interpretability of Generative EEG Microstate Segmentation}.\newline
  34th Irish Conference on Artificial Intelligence and Cognitive Science (AICS 2026), University College Dublin, 2026.
  \item \href{https://arxiv.org/abs/2605.10947}{Interpretable EEG Microstate Discovery via Variational Deep Embedding}.\newline
  XAI 2026, Late-breaking work and Doctoral Consortium track, Fortaleza, Brazil, 2026.
  \item \href{https://doi.org/10.1109/ICTAS59620.2024.10507142}{Machine Learning Models for Predicting Malaria in Nigerian Children Under Five}.\newline
  IEEE ICTAS 2024, 2024.
\end{itemize}

\end{document}
